\section{实验环境与评价指标}

实验在单张 GPU 环境下完成，软件框架采用 Python 与 PyTorch。为了更全面评估模型性能，本文选用分类准确率、参数量和单张图像推理时延作为主要评价指标。其中，分类准确率用于衡量模型识别能力，参数量和时延用于衡量模型的部署成本。

\section{对比实验结果}

\begin{table}[htbp]
  \centering
  \caption{不同方法在测试集上的性能比较}
  \label{tab:main-result}
  \begin{tabular}{lccc}
    \toprule
    方法 & 准确率 / \% & 参数量 / M & 推理时延 / ms \\
    \midrule
    Baseline-ResNet18 & 91.6 & 11.7 & 4.8 \\
    Baseline + 注意力模块 & 92.8 & 12.1 & 5.1 \\
    Baseline + 多尺度融合 & 93.1 & 12.4 & 5.3 \\
    本文方法 & 94.0 & 12.6 & 5.4 \\
    \bottomrule
  \end{tabular}
\end{table}

由表 \ref{tab:main-result} 可知，本文方法在准确率指标上优于各类对照模型，同时仅带来较小的参数量和时延增长，说明改进策略具备较好的精度效率平衡特性。

\section{消融实验}

为了分析不同模块的贡献，本文进一步开展消融实验，结果如表 \ref{tab:ablation} 所示。

\begin{table}[htbp]
  \centering
  \caption{模型消融实验结果}
  \label{tab:ablation}
  \begin{tabular}{cccc}
    \toprule
    通道注意力 & 多尺度融合 & 数据增强优化 & 准确率 / \% \\
    \midrule
    否 & 否 & 否 & 91.6 \\
    是 & 否 & 否 & 92.8 \\
    否 & 是 & 否 & 93.1 \\
    是 & 是 & 是 & 94.0 \\
    \bottomrule
  \end{tabular}
\end{table}

从表 \ref{tab:ablation} 可以看出，各模块均对最终结果有积极作用，其中多尺度融合对精度提升贡献略高于单独的通道注意力模块，而完整方案的性能最佳。

\section{结果分析}

通过对误分类样本进行观察，本文发现基线模型主要在以下三类场景中容易出现错误：

\begin{enumerate}
  \item 目标主体占比过小，导致深层网络难以保留足够的局部细节信息。
  \item 背景纹理较为复杂，模型对非目标区域存在错误关注。
  \item 同类别样本内部差异大，或者不同类别之间外观相似，造成边界模糊。
\end{enumerate}

针对上述问题，引入注意力机制能够强化判别区域响应，而多尺度特征融合则有助于保留纹理细节，因而两者在复杂场景下具有互补作用。整体来看，本文方法并非简单追求更深网络，而是通过结构优化改善特征利用效率，更符合实际部署需求。

\section{本章小结}

本章给出了实验环境、对比实验与消融实验结果，并从误分类样本分析角度解释了模型性能提升的原因。实验结果表明，本文所提出的改进策略在精度和效率之间实现了较好的平衡，验证了研究方案的有效性。
